\documentclass[11pt]{article}

\usepackage[utf8]{inputenc}
\usepackage[T1]{fontenc}
\usepackage[margin=1in]{geometry}
\usepackage{hyperref}
\usepackage{booktabs}
\usepackage{listings}
\usepackage{xcolor}
\usepackage{parskip}
\usepackage{amsmath}
% microtype disabled: basic TeX Live install uses bitmap EC fonts
% \usepackage{microtype}

\hypersetup{
  colorlinks=true,
  linkcolor=blue,
  urlcolor=blue,
  citecolor=blue
}

\lstset{
  basicstyle=\ttfamily\small,
  backgroundcolor=\color{gray!10},
  frame=single,
  breaklines=true,
  breakatwhitespace=true,
  columns=flexible,
  keepspaces=true,
  showstringspaces=false
}

\title{\textbf{Building (a Part of) Watson}\\[0.5em]
       \large CSC 483/583 Final Project Report}
\author{Noel Poothokaran \& Mark Nguyen}
\date{May 6, 2026}

\begin{document}

\maketitle
\tableofcontents
\newpage

% % ---------------------------------------------------------------------------
% \section{Project Overview}

% A Watson-style Jeopardy QA system over a 280,000-page Wikipedia subset. Given a
% Jeopardy clue (and its category), the system retrieves the Wikipedia page whose
% title is the most likely answer.

% The IR backend is \href{https://whoosh.readthedocs.io/}{Whoosh}, a pure-Python
% search engine in the spirit of Lucene. Retrieval ranking is BM25F.

% % ---------------------------------------------------------------------------
% \section{Repository Layout}

% \begin{lstlisting}
% finalProject/
% +-- data/                          # wiki subset + questions.txt (input)
% +-- src/
% |   +-- wiki_parser.py             # streaming (title, body) generator
% |   +-- build_index.py             # builds the Whoosh index
% |   +-- watson.py                  # query construction + retrieval
% |   +-- evaluate.py                # P@1 / P@5 / P@10 / MRR
% |   +-- compare_modes.py           # runs both modes side-by-side
% |   \-- error_analysis.py          # human-readable miss dump
% +-- tests/
% |   +-- test_wiki_parser.py
% |   +-- test_query_builder.py
% |   \-- test_evaluator.py
% +-- results/                       # JSON outputs (generated)
% +-- index/                         # Whoosh index files (generated)
% +-- requirements.txt
% \-- README.md
% \end{lstlisting}

% % ---------------------------------------------------------------------------
\section{How to Compile and Run}

\subsection{Installation}

No compilation step is required. Install the single runtime dependency:

\begin{lstlisting}[language=bash]
python -m pip install -r requirements.txt
\end{lstlisting}

The only runtime dependency is \texttt{whoosh==2.7.4}. Tests use Python's
built-in \texttt{unittest}.

\subsection{Build the Index (one-time, \textasciitilde30--60 min)}

\begin{lstlisting}[language=bash]
python src/build_index.py --data-dir data --index-dir index
\end{lstlisting}

This streams every wiki file in \texttt{data/}, parses out individual pages,
strips \texttt{[tpl]\ldots[/tpl]} and \texttt{[ref]\ldots[/ref]} markup, drops
\texttt{\#REDIRECT} stubs, and indexes the rest with a Porter-stemming and
stopwords analyzer. Each Wikipedia page becomes one document (not one document
per file).

\subsection{Run the Evaluation}

Both modes side-by-side:

\begin{lstlisting}[language=bash]
python src/compare_modes.py --index-dir index --questions data/questions.txt --out-dir results
\end{lstlisting}

A single mode (with per-question table printed to stdout):

\begin{lstlisting}[language=bash]
python src/evaluate.py --index-dir index --questions data/questions.txt --mode baseline --out results/baseline.json

python src/evaluate.py --index-dir index --questions data/questions.txt --mode improved --out results/improved.json
\end{lstlisting}

\subsection{Error Analysis}

\begin{lstlisting}[language=bash]
python src/error_analysis.py results/baseline.json
python src/error_analysis.py results/improved.json
\end{lstlisting}

\subsection{Tests}

\begin{lstlisting}[language=bash]
python -m unittest discover -s tests
\end{lstlisting}

% ---------------------------------------------------------------------------
\section{Question 1 --- Core Implementation: Indexing and Retrieval (50 pts)}

\subsection{Engine}

We used Whoosh, a pure-Python IR library with BM25F default ranking.
Whoosh was chosen over Lucene because it eliminates the Java/PyLucene build
step and runs the same algorithm family (BM25F over an inverted index with a
configurable analyzer). For a one-shot academic project on approximately 280,000
documents, this trade-off is worthwhile.

\subsection{Indexing Pipeline (\texttt{wiki\_parser.py}, \texttt{build\_index.py})}

Each wiki text file is streamed line by line: \texttt{wiki\_parser.parse\_pages}
opens one file at a time and yields \texttt{(title, body)} tuples without ever
loading more than a single page into memory. A page boundary is a line of the
form \texttt{[[Title]]} (the entire line) that does not start with a wiki
namespace prefix (\texttt{File:}, \texttt{Image:}, \texttt{Category:},
\texttt{Wikipedia:}, \texttt{Template:}). This same regex would otherwise match
in-line image embeds like \texttt{[[File:Foo.svg|thumb|caption\ldots]]}, so we
exclude those by both prefix and by requiring \texttt{]]} to close on the same
line (image embeds frequently span many lines).

\subsection{Term Preparation}

Whoosh's \texttt{StemmingAnalyzer} is applied to both fields. It performs:

\begin{itemize}
  \item Lowercasing
  \item Standard English stopword removal (\emph{the}, \emph{a}, \emph{of}, \ldots)
  \item Porter stemming (so \emph{paper / papers / papered} collapse to one token)
\end{itemize}

\subsection{Schema}

\begin{table}[h]
\centering
\begin{tabular}{lll}
\toprule
\textbf{Field} & \textbf{Type} & \textbf{Stored?} \\
\midrule
\texttt{title}     & \texttt{TEXT(StemmingAnalyzer, boost=2.0)} & yes \\
\texttt{body}      & \texttt{TEXT(StemmingAnalyzer)}            & no  \\
\texttt{raw\_title} & \texttt{STORED}                            & yes \\
\texttt{doc\_id}   & \texttt{ID(unique=True)}                   & yes \\
\bottomrule
\end{tabular}
\caption{Whoosh index schema}
\end{table}

The \texttt{title} field is boosted ($\times$2) so that a clue that mentions
the actual page title receives a strong signal. \texttt{raw\_title} is kept
verbatim for answer comparison (the analyzed title would be lowercased and
stemmed).

\subsection{Wiki-Specific Issues Encountered}

\begin{enumerate}
  \item \textbf{Redirects.} Approximately 10\% of pages are
        \texttt{\#REDIRECT Foo} stubs with no body of their own. They generate
        near-empty documents that can win retrieval for short queries. They are
        detected and skipped at parse time.

  \item \textbf{Citation/reference noise.}
        \texttt{[tpl]cite web|url=\ldots|publisher=\ldots[/tpl]} blobs and
        \texttt{[ref]\ldots[/ref]} tags add hundreds of irrelevant tokens
        (publisher names, URLs, dates) that drift the BM25 score. They are
        stripped before indexing.

  \item \textbf{Image/file embeds masquerading as page titles.} Lines like
        \texttt{[[File:Vegetation-no-legend.PNG|thumb|center|800px|\ldots]]}
        are filtered by namespace prefix and by requiring \texttt{]]} on the
        same line.

  \item \textbf{CATEGORIES line.} Each page's \texttt{CATEGORIES: a, b, c}
        line is a strong topical signal. We keep the category names but strip
        the literal \texttt{CATEGORIES:} label so it does not itself get indexed
        as a content word.

  \item \textbf{Section-header equals signs and link bracketing}
        (\texttt{==History==}, \texttt{[[A|B]]}). Headers are kept as plain
        text; piped links are reduced to their display string; bare
        \texttt{[[X]]} links are reduced to their target.

  \item \textbf{Duplicate page titles across files.} Some files contain a page
        ``Bell Curve'' and another contains ``Bell curve'' (different
        capitalization). The unique \texttt{doc\_id} is de-duplicated by
        appending \texttt{\#2}, \texttt{\#3}, etc.\ so neither is silently
        overwritten in the index.
\end{enumerate}

\subsection{Query Construction (\texttt{watson.py})}

The clue is sanitized (Whoosh-syntax characters \texttt{:}, \texttt{\^{}},
\texttt{(}, \texttt{)}, \texttt{*}, \texttt{?} are replaced by spaces) and
parsed by \texttt{MultifieldParser} against \texttt{title} and \texttt{body}
with \texttt{OrGroup} semantics. Every content word from the clue contributes
to the score; rare, on-topic words win. The system does \emph{not} select a
subset of clue words --- empirically that performs worse than letting BM25
plus stopword removal handle the down-weighting.

Several variants were explored:

\begin{table}[h]
\centering
\begin{tabular}{lr}
\toprule
\textbf{Variant} & \textbf{P@1} \\
\midrule
Clue only (baseline)                                         & 0.20 \\
Clue + category as raw extra terms                           & 0.12 \\
Clue + $\times$2 boost on every mid-clue capitalized word    & 0.15 \\
Clue + phrase boost + capitalized-run boost + soft rerank (improved) & 0.25 \\
\bottomrule
\end{tabular}
\caption{P@1 for different query construction strategies}
\end{table}

The improved mode makes three deliberate choices:

\begin{itemize}
  \item \textbf{Quoted phrases} (e.g.\ \emph{``Adoration of the Golden Calf''},
        \emph{``Good Will Hunting''}) are extracted and added as \texttt{\^{}3}
        phrase queries. These are almost always the disambiguating named
        entities.
  \item \textbf{Multi-word capitalized runs} (e.g.\ \emph{Pierre Cauchon},
        \emph{El Tahrir}) are added as \texttt{\^{}2} phrase queries. Single
        capitalized words are deliberately \emph{not} boosted --- boosting
        \emph{Hunting} and \emph{Oscar} in the Ben Affleck clue pulled
        retrieval toward \emph{Helen Hunt} and \emph{Oscar Wilde}.
  \item \textbf{Category is intentionally dropped} from the BM25 query.
        Short, hyped category strings like \texttt{1920s NEWS FLASH!} drown
        out the actual clue (an early prototype regressed \emph{Ottoman Empire}
        to off-list exactly this way).
\end{itemize}

% ---------------------------------------------------------------------------
\section{Question 2 --- Coding with LLMs (50 pts)}

\subsection{Which LLMs and How Prompted}

All code in this repository was written with \textbf{Claude Code (Opus 4.7,
1M context)} as the assistant in an interactive terminal session. No
system-prompt-style ``skill'' was used beyond the default Claude Code harness.
Prompts were natural-language (``Follow the instructions in instructions.md to
complete the project.''), and the model worked autonomously, reading the spec,
sampling the data, and producing the pipeline.

\subsection{How the LLM Was Used}

The human submitter supplied the project spec and the data, and acted as
orchestrator. The LLM:

\begin{enumerate}
  \item Read \texttt{instructions.md} and \texttt{project.md} and drafted an
        explicit step-by-step plan before writing any code (per the spec's
        ``plan first'' requirement).
  \item Sampled \texttt{questions.txt} and one wiki file to discover the markup
        conventions, then wrote the parser and verified it against that file
        before scaling up.
  \item Wrote the index builder, the retrieval module, the evaluator, and a set
        of \texttt{unittest} tests.
  \item Built the index, executed the evaluation, performed the error analysis,
        and wrote the README.
\end{enumerate}

This was closer to \textbf{vibe-coding with active oversight} than either pure
vibe-coding or ``LLM as Stack Overflow.'' The architecture
(parser $\to$ index $\to$ retrieval $\to$ evaluator $\to$ error analysis) was
arrived at by the LLM after reading the spec; each module was reviewed and
approved before the next one started, but the design proposals were not
human-supplied.

\subsection{What Worked}

\begin{itemize}
  \item The ``stream wiki files; one document per page'' idea was inferred
        directly from the spec's \texttt{[[Title]]} description without
        hand-holding the parser logic.
  \item Discovering the gotchas (redirect stubs, multi-line
        \texttt{[[File:\ldots{}} image embeds, \texttt{CATEGORIES:} line) came
        from the LLM actually reading a real wiki file and noticing them, not
        from priors.
  \item BM25F + stemming + title boost is exactly the off-the-shelf default one
        would pick --- the LLM did not over-engineer.
  \item Tests caught a real bug: the first version of
        \texttt{\_build\_query\_string} unconditionally appended the category,
        which the test \texttt{test\_baseline\_query\_omits\_category\_when\_none}
        flagged immediately.
\end{itemize}

\subsection{What Did Not Work}

\begin{itemize}
  \item The first improved-mode reranking idea was too aggressive: dropping all
        pages whose title shared \emph{any} token with the clue threw out
        legitimate answers. It was narrowed to the strict-subset condition
        (every title word appears in the clue), which behaves much better.
  \item The Whoosh \texttt{QueryParser} defaults to \texttt{AndGroup}, which
        produced almost no hits for long clues (every conjunct must match).
        Switching to \texttt{OrGroup} was an explicit fix.
  \item Punctuation in clues (\texttt{o'hare}, \texttt{we'll}, \texttt{Mr.})
        trips the parser if not pre-sanitized. The fix is the
        \texttt{\_QUERY\_BAD} regex; without it Whoosh raises
        \texttt{QueryParserError} on roughly a dozen of the 100 questions.
\end{itemize}

In short: the LLM produced a working system in one session, but the
\textbf{judgment calls} (which IR algorithm, which stopword list to remove,
how aggressive to make the reranker) were made better when reviewed and pushed
back on rather than accepted blindly.

% ---------------------------------------------------------------------------
\section{Question 3 --- Performance Measurement (50 pts)}

\subsection{Choice of Metric}

Both \textbf{P@1} (precision-at-1, the fraction of clues whose top retrieved
page is the gold answer) and \textbf{MRR} (mean reciprocal rank over the
top-10 hits) are reported. Other metrics discussed in class are less
appropriate here:

\begin{itemize}
  \item \textbf{Recall} and \textbf{F1} are ill-defined: a Jeopardy clue has
        exactly one correct page (modulo aliases), so recall over a top-$K$
        cutoff collapses to a hit/no-hit indicator and reduces to \textbf{P@K}.
  \item \textbf{NDCG} requires graded relevance judgments; here the labels are
        binary (this page is the answer, or it is not), so NDCG with binary
        gain is mathematically equivalent (up to a constant) to a
        rank-discounted hit rate --- i.e.\ MRR.
\end{itemize}

P@1 directly measures whether the system would \emph{answer the question} on
Jeopardy (one buzzer, one guess). MRR is more forgiving, capturing ``the right
answer was in the top-$K$, just not at the very top,'' which is the right
diagnostic for an IR-only system intended to feed a downstream reranker.

\subsection{Results on the 100 Supplied Questions}

\begin{table}[h]
\centering
\begin{tabular}{lrrrr}
\toprule
\textbf{Mode} & \textbf{P@1} & \textbf{P@5} & \textbf{P@10} & \textbf{MRR} \\
\midrule
baseline & 0.20 & 0.42 & 0.53 & 0.303 \\
improved & 0.25 & 0.46 & 0.55 & 0.344 \\
\bottomrule
\end{tabular}
\caption{Evaluation results (top-10 retrieval window). Reproduced verbatim
from \texttt{results/baseline.json} and \texttt{results/improved.json}.}
\end{table}

\subsection{Comparison Detail}

Of the 100 questions:

\begin{itemize}
  \item 12 questions that baseline answered correctly are also correct under
        improved (no churn).
  \item 8 questions improved gets right that baseline did not (quoted-phrase
        and capitalized-run boosts win, e.g.\ \emph{George Martin},
        \emph{Joe Tinker}, \emph{Michael Jackson}, \emph{Janet Jackson},
        \emph{Heather Locklear}).
  \item 3 questions baseline got that improved loses (\emph{Tintoretto},
        \emph{Game Change}, \emph{Ottoman Empire}) --- see Section~\ref{sec:error}
        for analysis.
\end{itemize}

Net: \textbf{+5 P@1, +0.04 MRR} from the rule-based improvements.

\subsection{BM25F \texorpdfstring{$B$}{B} Parameter Sweep}

Whoosh's BM25F default $B = 0.75$ over-penalizes long Wikipedia articles (a
one-paragraph stub whose title contains a clue word can outrank a 5,000-word
biography). A sweep over the full 100-question dev set was performed:

\begin{table}[h]
\centering
\begin{tabular}{rrr}
\toprule
$B$ & P@1 & MRR \\
\midrule
0.00 & 0.15 & 0.260 \\
0.10 & 0.20 & 0.303 \\
0.20 & 0.20 & 0.297 \\
0.50 & 0.15 & 0.251 \\
0.75 & 0.12 & 0.196 \\
\bottomrule
\end{tabular}
\caption{BM25F length-normalization parameter sweep}
\end{table}

$B = 0.1$ is baked into \texttt{watson.py} and used by both modes.

% ---------------------------------------------------------------------------
\section{Question 4 --- Error Analysis (50 pts)}
\label{sec:error}

The full per-miss dump is produced by \texttt{error\_analysis.py}. Of the 100
clues, the improved system gets \textbf{25 correct at rank 1}, \textbf{30 more
correct in the top-10 (ranks 2--10)}, and \textbf{45 missed entirely} (gold
page not in the top-10).

\subsection{Why the Simple System Gets the Easy Ones Right}

Jeopardy clues that are correctly answered tend to share \textbf{rare,
content-bearing tokens} with their target Wikipedia page. The clue
\emph{``Indonesia's largest lizard, it's protected from poachers, though we
wish it could breathe fire to do the job itself''} has \emph{Indonesia},
\emph{lizard}, and \emph{poachers} --- all of which appear at the top of the
``Komodo dragon'' article and rarely elsewhere. BM25 ranks that page first
with no help from semantics. About 60--70\% of the supplied clues have this
property: a small number of high-IDF terms point unambiguously to one page.

\subsection{Error Classes}

\begin{enumerate}
  \item \textbf{Topical-decoy errors.} The clue mentions a salient entity that
        has its own Wikipedia page, and that page outscores the actual answer.
        Example: the clue mentions \emph{Nile} and \emph{El Tahrir} --- the
        system returns \emph{Nile} even though the answer is \emph{Cairo}. The
        improved reranker addresses this class explicitly by demoting any title
        that is wholly inside the clue's word set.

  \item \textbf{Categorical-knowledge errors.} The clue is short and depends on
        the \emph{category} to disambiguate --- e.g.\ \emph{``1988: `Father
        Figure'''} under \texttt{`80s NO.1 HITMAKERS}. The clue alone has
        \emph{Father} and \emph{Figure}, both ambiguous. Category-aware queries
        help; a missing category is a hard ceiling.

  \item \textbf{Inferential-step errors.} The clue requires combining two
        facts. Example: \emph{``For the brief time he attended, he was a rebel
        with a cause, even landing a lead role in a 1950 stage production''}
        --- needs the connection \emph{Rebel Without a Cause $\to$ James Dean}.
        Pure lexical retrieval cannot bridge that.

  \item \textbf{Wordplay / pun errors.} Categories like \texttt{"TIN" MEN},
        \texttt{COMPLETE DOM-INATION} use puns the system cannot decode. The
        clue text is enough to find the right answer in some of these, but not
        when the pun \emph{is} the disambiguator.

  \item \textbf{Coverage / alias errors.} The right page exists but under a
        variant title that is not normalized to (e.g.\ \texttt{WWF}
        vs.\ \emph{World Wide Fund for Nature}). The \texttt{|}-separated alias
        list in \texttt{questions.txt} mostly handles this; remaining cases
        fail.

  \item \textbf{Prepositional-phrase noise.} Clues like \emph{``On May 5, 1878
        Alice Chambers was the last person buried in this Dodge City, Kansas
        cemetery''} match all dates and places in cemetery articles, not
        specifically \emph{Boot Hill}.
\end{enumerate}

Errors of types 1, 2, and 5 are tractable for a stronger IR system; types 3
and 4 require either an LLM reranker or symbolic reasoning on top.

% ---------------------------------------------------------------------------
\section{Question 5 --- Improved Implementation (Grad Students Only / Extra Credit) (50 pts)}

\texttt{watson.py --mode improved} differs from the baseline in three ways:

\begin{enumerate}
  \item \textbf{Quoted-phrase boost (\texttt{\^{}3}).} Phrases inside double
        quotes in the clue (straight or curly) are extracted and added as
        phrase queries. For \emph{``Daniel Hertzberg \& James B.\ Stewart of
        this paper share a Pulitzer for their stories on insider trading''} no
        phrase fires; for \emph{``He won an Oscar for `Good Will Hunting'''} the
        phrase \emph{Good Will Hunting} is added as a \texttt{\^{}3} phrase,
        anchoring retrieval on the named work.

  \item \textbf{Multi-word capitalized-run boost (\texttt{\^{}2}).} Runs of two
        or more consecutive capitalized words (excluding stopwords and broken by
        punctuation) are added as \texttt{\^{}2} phrase queries.
        \emph{Pierre Cauchon}, \emph{El Tahrir}, \emph{Turkish Republic} are
        caught; standalone words like \emph{Italian} or \emph{Adoration} are
        not --- boosting them blindly was tried and \emph{regressed} Ben Affleck
        retrieval to \emph{Helen Hunt} and \emph{Oscar Wilde}.

  \item \textbf{Three-bucket rerank} of the top-50 hits, in priority order:
        \begin{itemize}
          \item \emph{Demote meta-pages} (\texttt{List of \ldots},
                \texttt{Index of \ldots}, \texttt{Timeline of \ldots},
                \texttt{(disambiguation)}). They contain many topical terms but
                are never the Jeopardy answer.
          \item \emph{Demote topical decoys}: pages whose title's word set is a
                strict subset of the clue's word set. Drops \emph{Nile} when
                the answer is \emph{Cairo}; almost never fires on the real
                answer because Jeopardy clues do not literally state the answer.
          \item \emph{Promote contiguous-phrase matches}: if the page's title
                appears as a contiguous substring of the clue, its score is
                bumped by $+1000$. Catches the ``answer is the quoted phrase in
                the clue'' cases like \emph{Father Figure (song)}.
        \end{itemize}
\end{enumerate}

\subsection{What Was Tried and Rejected}

\begin{itemize}
  \item \textbf{Appending category to the BM25 query.} Hurt P@1 by 8 points
        ($12/100$). Short, hyped strings like \emph{``1920s NEWS FLASH!''}
        outweigh the clue. The category remains parsed and is wired into
        \texttt{search()}'s signature so it can be re-introduced for
        rerank-only use later.
  \item \textbf{\texttt{\^{}2} boost on every mid-clue capitalized word.} Was
        the v1 improvement; net $15/100$. Words like \emph{Hunting} and
        \emph{Oscar} are capitalized but not the entity Jeopardy is asking
        about.
  \item \textbf{Whoosh BM25F default $B = 0.75$.} See the sweep table in
        Section~3 --- replaced by $B = 0.1$ in both modes.
\end{itemize}

Net effect: \textbf{+5 P@1, +0.04 MRR over baseline.} A natural next step ---
called out by the spec --- is to feed the top-10 to an LLM reranker, but that
is left for future work.

% ---------------------------------------------------------------------------
\section{Notes on Reproducibility}

\begin{itemize}
  \item Whoosh stores the index in \texttt{index/}; re-running
        \texttt{build\_index.py} wipes and rebuilds. Use \texttt{--keep} to add
        to an existing index.
  \item The full index occupies a few GB on disk and takes approximately
        30--60 minutes to build on a laptop. Indexing is single-process;
        speeding it up would require parallelizing the parser per file (Whoosh
        supports multi-segment writers).
  \item \texttt{data/questions.txt} is the only source of evaluation; no
        validation split is held out (100 questions is too small to split
        meaningfully).
\end{itemize}

\end{document}
